\chapter{أوجه الاختلاف بين الأحاديث: الجمع والتخصيص}

\section{حكم قتل النساء والصبيان في البيات (الغارات الليلية)}
يواصل الإمام الشافعي رحمه الله مناقشة أوجه الاختلاف الظاهري بين الأحاديث، ومن ذلك:
حديث الصعب بن جثامة رضي الله عنه: أن النبي ﷺ سُئل عن أهل الدار من المشركين يُبَيَّتون (أي يُهجم عليهم ليلاً) فيُصاب من نسائهم وذراريهم، فقال: «هم منهم» وفي رواية «هم من آبائهم». وظاهره إباحة إصابتهم.
ويقابله حديث ابن كعب بن مالك: أن النبي ﷺ لما بعث إلى ابن أبي الحُقَيق نهى عن قتل النساء والولدان. وظاهره التحريم.

\textbf{طريقة الجمع أو الترجيح:}
ذهب سفيان بن عيينة إلى أن حديث النهي ناسخ لحديث الإباحة (هم منهم). لكن الشافعي لم يقنع بهذا، وناقشه نقاشاً علمياً وتاريخياً؛ فبين أن حديث الصعب كان في عمرة النبي ﷺ المتأخرة، ويقيناً كان بعد مقتل ابن أبي الحقيق، فلا يمكن أن يكون المتقدم ناسخاً للمتأخر. 
ثم جمع الشافعي بينهما بفقه دقيق: أن النهي عن قتل النساء والولدان يُحمل على (التعمد والقصد) وهم متميزون بعيدون عن المقاتلة؛ لأنهم لم يبلغوا كفراً يُقاتلون عليه، ولا معنى لقتالهم. أما حديث «هم منهم» فيُحمل على البيات (الغارة الليلية) أو اختلاطهم بالمقاتلين بحيث لا يمكن تمييزهم، فإذا أُصيبوا خطأً فلا دية فيهم ولا كفارة ولا قود، لأنهم كفار في دار كفر، وليس لهم حكم الإيمان الذي يعصم الدم.

\subsection{دلالة من القرآن على سقوط الدية والكفارة في دار الكفر}
استدل الشافعي على سقوط الكفارة والدية عنهم بآيات قتل الخطأ في سورة النساء؛ فالمؤمن في دار الإيمان (أو المعاهد) فيه الكفارة والدية ﴿وَدِيَةٌ مُّسَلَّمَةٌ إِلَىٰ أَهْلِهِ وَتَحْرِيرُ رَقَبَةٍ﴾. أما المؤمن الذي في دار الكفر ففيه الكفارة دون الدية ﴿فَإِن كَانَ مِن قَوْمٍ عَدُوٍّ لَّكُمْ وَهُوَ مُؤْمِنٌ فَتَحْرِيرُ رَقَبَةٍ﴾. فإذا كان المؤمن تسقط ديته في دار الكفر، فمن باب أولى أن يسقط القود والدية والكفارة عن نساء المشركين وصبيانهم إذا أُصيبوا خطأً في البيات.

\section{غسل الجمعة بين الوجوب والاستحباب}
ومن الأحاديث التي اختلف في فقهها:
حديث أبي سعيد الخدري: «غسل الجمعة واجب على كل محتلم». وحديث ابن عمر: «من جاء منكم الجمعة فليغتسل». ظاهرها الوجوب الذي لا تصح الصلاة إلا به.

\textbf{توجيه الشافعي وتخصيصه:}
استدل الشافعي بحديث دخول عثمان بن عفان رضي الله عنه المسجد يوم الجمعة وعمر يخطب، فعاتبه عمر لعدم اغتساله واكتفائه بالوضوء، فلو كان الغسل واجباً شرطياً (كركن أو شرط للصلاة) لأمره عمر بالخروج ليغتسل، ولما واصل عثمان صلاته بدونه. فدل إقرار عمر وعثمان والصحابة على أن الغسل هنا (واجب في الاختيار والنظافة) وليس شرطاً في صحة الصلاة.
ويؤيده حديث: «من توضأ يوم الجمعة فبها ونعمت، ومن اغتسل فالغسل أفضل». فالأوامر تُحمل على الاستحباب وطلب الأفضلية، لا على الركنية.

\section{الخطبة على خطبة الأخ والبيع على بيعه}
قال ﷺ: «لا يخطب أحدكم على خطبة أخيه». ظاهره التحريم المطلق بمجرد الخطبة.
لكن الشافعي بين أن هذا النهي مصروف لمعنى خاص؛ واستدل بحديث فاطمة بنت قيس رضي الله عنها، حين طلقها زوجها، فجاءت تستشير النبي ﷺ وأخبرته أن معاوية وأبا جهم خطباها، فقال ﷺ: «أما أبو جهم فلا يضع عصاه عن عاتقه، وأما معاوية فصعلوك لا مال له، انكحي أسامة بن زيد».
فالنبي ﷺ خطبها لأسامة بعد أن علم بخطبة معاوية وأبي جهم، فدل ذلك على أن النهي (لا يخطب أحدكم على خطبة أخيه) إنما يكون بعد (الركون والقبول) من المرأة ووليها. أما قبل الموافقة والقبول، فالحال واحدة، ويجوز للرجل أن يتقدم لخطبتها كما فعل النبي ﷺ لأسامة.

\subsection{لا يبيع الرجل على بيع أخيه}
وهو مقاس على ما قبله؛ لقوله ﷺ: «المتبايعان بالخيار ما لم يتفرقا»، فنهيه عن البيع على بيع أخيه إنما يكون قبل التفرق (في وقت الخيار)، فإذا جاء مشترٍ آخر ليفسد البيعة بزيادة السعر، فهذا هو المحرم. أما إذا تم البيع ولزم العقد وتفرقا، فلا يضره أن يبيع أو يشتري لأن البيع قد لزم ولا يملك فسخه.

\section{الصلاة في أوقات النهي (الفرائض وذوات الأسباب)}
نهى النبي ﷺ عن الصلاة بعد العصر حتى تغرب الشمس، وبعد الصبح حتى تطلع الشمس. 
واحتمل النهي أن يكون عاماً لكل الصلوات، واحتمل أن يكون مخصصاً للنوافل. فبحث الشافعي عن الدلائل، فوجد أن النبي ﷺ قال: «من أدرك ركعة من الصبح قبل أن تطلع الشمس فقد أدرك الصبح»، وهذا يقتضي أنه سيُكمل صلاته أثناء طلوع الشمس وهو وقت نهي. وقال: «من نسي صلاة فليصلها إذا ذكرها»، وجعل وقت الذكر هو وقتها ولو كان في وقت نهي.
وقال في الطواف: «يا بني عبد مناف لا تمنعوا أحداً طاف بهذا البيت وصلى أية ساعة شاء من ليل أو نهار».

فدل هذا كله على أن النهي مخصص بـ (النوافل المطلقة) التي لا سبب لها. أما الصلوات المفروضة، والصلوات التي لها سبب (كسنة الطواف، وركعتي الكسوف، والجنائز) فتُصلى في أي وقت، ولا تدخل في النهي. 
وقد طاف ابن عمر وابن عباس والحسن والحسين وصلوا ركعتي الطواف بعد العصر وبعد الصبح، وأما تأخير عمر بن الخطاب لركعتي الطواف حتى طلعت الشمس، فلأنه حفظ النهي المجرد ولم تبلغه الرخصة، ومن علم حجة على من لم يعلم.

\section{وجوب التسليم للسنة إذا ثبتت}
يختم الشافعي هذا الباب بكلمة عظيمة في منهج الاستدلال: «وإذا ثبت عن رسول الله ﷺ الشيء، فهو اللازم لجميع من عرفه، لا يقويه ولا يوهنه شيء غيره، بل الفرض الذي على الناس اتباعه... ولم يجعل الله لأحد معه أمراً يخالف أمره». 
فالسنة حاكمة بنفسها، ولا تحتاج لعمل الناس لتقويتها، ولا يضعفها ترك بعض الناس لها.